El sistema híbrido desarrollado en el \cref{ch:cap4} demostró capacidad predictiva para estimar emisiones de \ce{CO2} y detectar eventos de anaerobiosis mediante \ce{CH4} con métricas de desempeño aceptables a escala de laboratorio. Sin embargo, la validación técnica del motor DEB+GP constituye solo una condición necesaria, no suficiente, para la gestión ambiental operativa. Mientras el Capítulo 4 resolvió el problema de estados parcialmente observables —estimando biomasa larval indirectamente y calibrando parámetros bioenergéticos—, su arquitectura permanece desconectada del ciclo de toma de decisiones del productor agroindustrial. Un modelo predictivo con MAPE del \num{8.3}~\% y un detector de \ce{CH4} con Recall del \num{100}~\% generan conocimiento técnico, pero no instrucción operativa. La pregunta que este capítulo aborda es, por tanto, cómo traducir esas predicciones en acciones de manejo concretas, trazables y verificables en contextos reales de bioconversión.

La brecha entre predicción y acción no es meramente técnica; es epistémica y organizacional. Como se discutió en el \cref{ch:marco_teorico}, la gobernanza ambiental efectiva requiere que los datos sean no solo precisos, sino \textit{actionables} \parencite{adadiPeekingBlackBoxSurvey2018}. Un operario de una planta de bioconversión en el Valle del Cauca no requiere la curva de respiración metabólica de las larvas, sino una indicación clara sobre si el lecho debe voltearse, si la humedad debe ajustarse o si el proceso acumula créditos de carbono verificables. Esta traducción exige una arquitectura de soporte a la decisión que opere como interfaz semántica entre el modelo matemático y el conocimiento empírico del productor \parencite{kamilarisReviewPracticeBig2017}.

Para cerrar esta brecha, este capítulo propone una plataforma de tres capas que extiende la arquitectura híbrida del Capítulo 4 hacia la aplicación operativa:

\begin{enumerate}
    \item \textbf{Capa de percepción y edge computing:} nodos IoT basados en ESP32 con sensores NDIR de \ce{CO2} y \ce{CH4}, complementados con sensores de temperatura y humedad, desplegados en condiciones de autonomía rural donde la conectividad es intermitente \parencite{duobieneDevelopmentWirelessSensor2022}. El procesamiento en el borde permite ejecutar algoritmos de filtrado y detección de anomalías localmente, garantizando continuidad funcional incluso ante fallos de red \parencite{rajRealTimeEstimation2023}.
    
    \item \textbf{Capa de procesamiento híbrido:} motor integrado DEB+GP, respaldado por una base de datos relacional (PostgreSQL) y un backend que orquesta la ejecución secuencial del estimador de biomasa, la EDO calibrada y el detector de residuos. Esta capa materializa los resultados del \cref{ch:cap4} en un pipeline automatizado de ingesta, inferencia y almacenamiento histórico.
    
    \item \textbf{Capa de aplicación:} dashboard web y sistema de alertas que traduce el inventario dinámico de carbono ($\Delta_{\text{Total}} = +\num{5.1}$~g~C, \num{3.6}~\% del total medido, según validación del Capítulo 4) en indicadores semafóricos de riesgo ambiental, junto con reportes estructurados para protocolos MRV \parencite{verraVCSStandardV472024}.
\end{enumerate}

La elección de una arquitectura descentralizada responde a tres restricciones contextuales del sur global. Primero, la conectividad rural inestable exige que la inteligencia del sistema no dependa exclusivamente de la nube \parencite{herrinCollateralDataQuality2021}. Segundo, el costo marginal del sistema debe ser bajo respecto al valor del producto (biomasa y frass) para garantizar adopción por pequeños productores. Tercero, la soberanía de datos —principio epistemológico defendido en el \cref{sec:gobernanza_datos}— requiere que la información ambiental permanezca bajo control local, evitando extractivismo de datos hacia plataformas centralizadas externas \parencite{fawzyStrategiesMitigationClimate2020}.

El presente capítulo desarrolla, en consecuencia, la especificación funcional y los protocolos operativos que habilitan la transición del sistema híbrido validado en el \cref{ch:cap4} a una herramienta de gestión ambiental integral. La \cref{sec:cap5_arquitectura} detalla la arquitectura tecnológica de las tres capas, especificando requerimientos de hardware y software sin incluir código de implementación, el cual se remite al Anexo~D. La \cref{sec:cap5_semaforo} define el semáforo de riesgo ambiental y los estados operativos asociados (verde, ámbar, rojo), vinculando cada nivel con acciones de mitigación concretas. La \cref{sec:cap5_protocolo} formaliza el protocolo de respuesta ante eventos de anaerobiosis, estableciendo tiempos de respuesta esperados y efectividad de las intervenciones. La \cref{sec:cap5_mrv} aborda la generación automatizada de reportes de carbono alineados con estándares IPCC Tier~3 y Verra VCS, incluyendo estructura de trazabilidad y sellado de datos. La \cref{sec:cap5_validacion} presenta métricas de utilidad operativa y pruebas de campo orientadas a operadores no expertos. Finalmente, la \cref{sec:cap5_conclusiones} sintetiza los aportes del capítulo y establece el puente hacia las conclusiones generales de la tesis.

Es importante precisar que este capítulo no pretende sustituir el juicio técnico del operador, sino estructurarlo. Como se argumentó en el \cref{sec:limites_reduccionismo}, el sistema actúa como traductor que hace visible lo invisible (las emisiones gaseosas), pero la decisión final de manejo permanece anclada en el contexto ecológico y productivo local. La plataforma de soporte a la decisión es, por diseño, una herramienta descriptiva y normativa de segundo orden: no impone reglas desde fuera, sino que externaliza las consecuencias ambientales cuantificadas de las decisiones operativas ya existentes.
